% figures/riscv-programming-spectrum_gradient.tex
% VARIANT (_gradient, ch1 house style) — not \input by main.tex; promote manually.
% Mirrors docs/thesis/figures/ch1-introduction/acceleration-spectrum.tex:
% white boxes chained left->right, a single shaded gradient bar on top encoding
% the control<->optimisation trade-off, and a «Παρούσα εργασία» figpill on the
% method this thesis actually uses (C Intrinsics for the crypto kernels).
% NOTE: move the pill to another box if the «current work» method differs.
\resizebox{\textwidth}{!}{%
\tikzsetnextfilename{ch3-riscv-riscv-programming-spectrum_gradient}
\begin{tikzpicture}[
    node distance=0.8cm,
    box/.style={
      draw=figrule, rounded corners=3pt, minimum width=3cm, minimum height=1.4cm,
      align=center, figT/label, fill=white
    },
    highlight/.style={box, draw=figaccent, figLW/emph, fill=figaccentfill},
    endlabel/.style={figT/annot},
  ]

  % Boxes (left = most control, right = most compiler optimisation)
  \node[box]                              (asm)    {Pure Assembly\\[-2pt]\footnotesize (\texttt{.S})};
  \node[box, right=0.8cm of asm]          (inline) {Inline Assembly\\[-2pt]\footnotesize (\texttt{\_\_asm\_\_})};
  \node[highlight, right=0.8cm of inline] (intr)   {C Intrinsics\\[-2pt]\footnotesize (\texttt{\_\_riscv\_*})};
  \node[box, right=0.8cm of intr]         (auto)   {Auto-Lowering\\[-2pt]\footnotesize (Plain C)};

  % Gradient spectrum bar: control[vermillion] -> optimisation[green]
  \shade[left color=figthreat!42, right color=figsolution!35, rounded corners=2pt]
    ([yshift=0.70cm]asm.north west) rectangle ([yshift=1.18cm]auto.north east);
  \node[endlabel, anchor=west] at ([xshift=2pt,yshift=0.94cm]asm.north west)
    {$\leftarrow$ Μεγαλύτερος έλεγχος};
  \node[endlabel, anchor=east] at ([xshift=-2pt,yshift=0.94cm]auto.north east)
    {Βελτιστοποίηση compiler $\rightarrow$};

  % «You are here» pill on the C Intrinsics box
  \node[figpill, below=0.55cm of intr] (here) {Παρούσα εργασία};
  \draw[draw=figaccent, figLW/base] (intr.south) -- (here.north);

\end{tikzpicture}%
}% end resizebox
